\begin{proposicion}[Bound by support and number of cycles]\label{prop:soportes-llegada}
Let \(U,v,c_s,d_i\) be as in Lemmas~\ref{lem:rango-llegada} and~\ref{lem:tipos-llegada}. Let \(\widehat
e_U\) be the logarithm of the sum of the weights of Proposition~\ref{prop:menores-llegada} on each face,
minus the change \(\Lambda\) in the predictor upon activating the pairs of that face, with the same
baseline normalization. Suppose that every site has a diagonal neighbor in \(U\). For \(v\le8\log n\),
\[
 |\Delta_U\widehat e_U|\le C^v(Cv)^{c_s+1}j^{-p}\prod_{i\in U}d_i,\qquad
 p=\begin{cases}v+c_s+1&\text{forest},\\v+c_s&\text{one cycle},\\
 v+c_s-1&\text{two or more independent cycles}.\end{cases}
\]
For \(v\le4\) the same orders hold for the physical error \(e\).
\end{proposicion}
\begin{proof}
Introduce a size variable \(s\) and independent activities \(\zeta_i\) after extracting the affine field
common to the faces. At the physical value, \(s=1/j\) and \(|\zeta_i|\le C/j\). A fiber of size \(D_f\)
satisfies \(\sum_{i\in f}d_i^{-1}\le1\). With \(\|z\|_*=\max_i d_i|z_i|\), the change in fiber load is
\(\le4\|z\|_*\) per neighborhood, and the measures for the exterior balancing vary by \(O(v|s|)\). Its
centered inverse and the propagators remain bounded as in Proposition~\ref{prop:menores-llegada}. On
\(|s|<\kappa/v,\ |\zeta_i|<\kappa'/d_i\), the derivative of the local odds equation in the \(*\) norm is
at most \(C\kappa'\). The normalized Hessian of the prefactor is \(I+E\), with \(\|E\|_\infty\le
C\kappa'\). The trace series bounds its logarithm by \(v\sum_{k\ge1}\|E\|_\infty^k/k\). In the
approximate count, each site enters by the empty-plus-occupied identity. The occupied-to-empty ratio
carries the factor \(\zeta_i\). When its neighbors are forbidden, the variations sum to \(O(\sum_{l\sim
i}|\zeta_l|)=O(\kappa')\). In an atom the deletions enter the measure variable \(\theta=sT_U\), of fixed
radius in total variation. The derivatives of order \(h\) of the exterior field, with its pole \(1/s\),
give mixed differences \(\le h! C^h|s|^{h-1}\) upon integration in the \(h\) deleted directions. Summing
them over subsets of \(h-1\) additional sites costs a geometric series with polynomial factor \(h\) if
\(Cv|s|<1\). If the occupied-to-empty ratio is at most \(K|\zeta_i|<1/2\) in smaller volumes, the
successive complex marginals have absolute value at most one. Their product controls every indicator of
prescribed sites. The change in the exterior field upon occupying \(i\) costs at most
\(\prod_D(1+|e^{a_D}-1|)\le e^{\sum_D|a_D|}\), where \(a_D\) are the mixed differences. Removing the
neighbors costs \(\exp(2K\sum_{l\sim i}|\zeta_l|)\). First fixing \(K\) for the exterior field and then
choosing \(\kappa'\) small closes the induction. Thus the partition function ratios are
\(1+O(|\zeta_i|)\), and the logarithm is \(O(v)\) per face. The saddle and cofactor factors agree on two
faces when \(\zeta_i=0\). Hence the alternating sum is holomorphic, bounded by \(C^v\), and divisible by
\(\prod_i\zeta_i\).

It remains to count degrees in \(s\). Upon deleting the marked fibers, an additional component mass
enters as \(s y_f\), where \(y_f\) sums the occupations of that fiber. The action has an outer power
\(s^{-1}\), and the initial power in the prefactor is \(s^0\). Thus an ordinary interaction of cost
\(\ell\) has at most \(\ell+1\) mean factors in the action and \(\ell\) in the prefactor. The term using
both marked measures in the second moment of the kernel has common amplitude \(s^2 z^2/2\). For the
masses \(y_i,y_h\), write \(F(y)=(1-y)^2-1,\ L(y)=-y\). Their polynomials are \(F_iF_h\) in the action
and \(F_iF_h-L_iF_h-F_iL_h\) in the prefactor. Their multiplier at cost \(\ell\) has degree
\(\le\ell-2\). Terms of positive excess already carry \(s\), so they can meet at most \(\ell-1\) groups
at that cost.

The logarithm turns the factorization of contributions on separate groups into a sum. For a connected
monomial of total cost \(w\) whose interaction \(a\) costs \(\ell_a\) and meets \(h_a\) groups, Euler's
formula gives the number of groups met as
\[
 1+w-\sum_a(\ell_a+1-h_a)-L_I ,
\]
where \(L_I\) counts the cycles of its bipartite incidence graph. Each deficit is nonnegative. If a group
is missing, the projection vanishes. The initial cost is at least \(c_s-1\).

In a forest, the branches of capacity-one factors can be marginalized exactly. Center the sites at the
mean of the tilted saddle. With \(p_0\) centered legs, \(u\) mean factors, \(h\) cumulant blocks, and
\(n_R,n_A\) ordinary action and prefactor vertices, there are \(L=p_0-n_R-n_A-h+1\) cycles. Since
\(p_0+u\le w+n_R\), at most \(u+h\le w+1-n_A-L\) groups are met. The centered blocks and action vertices
have degree at least two. With a single prefactor vertex, a tree is impossible because there is only one
possible leaf. If \(n_A=0,L=1\), every degree is two. The series of these cycles is
\(-\tfrac12\log\det(I-VR'')\), with \(V\) the covariance by group. The prefactor adds the same
half-logarithm with the opposite sign. Decorations of the link field of cost \(t\) add at most \(t\)
groups, since they differentiate the action of degree \(\le t+1\). With an exceptional interaction of
deficit one, meeting \(w\) groups requires a tree. When a branch is pruned, the means of its leaves
multiply the derivative in the link field, both in the count and at the saddle. The cross blocks of its
Hessian have zero trace. At the root the two fibers belong to different groups. For \(m_i=\mathbb Ey_i\),
\(v_i=m_i(1-m_i)\), the stationary contribution after taking the trace is
\(2F_iF_h+(m_i+v_i)F_h+(m_h+v_h)F_i=0\), since \(F(m_i)=-m_i-v_i\). On the binary atoms the sum of the
two polynomials vanishes. The identity persists under differentiation in the fields. Thus the forest
meets at most \(w-1\) groups.

In the unicyclic case, prune the outer branches at the critical cost \(c_s-1\). On the cycle the binary
means \(x_i\) give the transition from \(i\) to its successor \(h\), namely
\(\left(\begin{smallmatrix}1-t&t\\1&0\end{smallmatrix}\right)\), \(t=x_h/(1-x_i)\). The trace of the
product is \(1+\theta\), with \(\theta=(-1)^{|\mathcal C|}\prod_{i\in\mathcal C}x_i/(1-x_i)\). For the
cofactor, let \(E\) be the matrix copying sites to incidences and \(C_f\) the categorical covariance on
the fiber divided rowwise by its diagonal of variances. With \(C=\bigoplus_f C_f\), the determinant lemma
transforms the local cofactor \(\det C\,\det[I-E^{\mathsf T}E+E^{\mathsf T}C^{-1}E]\) into
\(\det[I-(C-I)(EE^{\mathsf T}-I)]\). A circuit changes site and fiber at each step, so only the two
orientations of the cycle survive, each of weight \(\theta\). The determinant equals \((1-\theta)^2\).
The term remaining in the error is \(\log(1-\theta^2)\). The quantity \(x_i\) is divisible by
\(\zeta_i\), and \(\zeta_i\partial_{\zeta_i}\) preserves that double multiplicity.

These coefficient identities persist when excesses are truncated at \(q=4v+12\). The moment recurrence
gives order \(\lceil k/2\rceil\) for \(\mu_k\), and even after the prefactor \(s^{-1}\) the change starts
at order \(\ge q-1>c_s+1\). The holomorphic identity extends the vanishing to the stated radii. If
\(R=\kappa/(2v)\) and \(r_i=\kappa'/(2d_i)\), Cauchy's estimate in \(s\) and Schwarz's lemma in each
field bound the tail starting at \(w\) by \(C^v\prod_i(|\zeta_i|/r_i)(|s|/R)^w/(1-|s|/R)\). The initial
costs are \(c_s+1,c_s,c_s-1\), respectively, except for the unicyclic coefficient of cost \(c_s-1\),
which gains \(\prod_{i\in\mathcal C}|\zeta_i|/r_i
\le Cv^2/j\). At the physical point this gives the stated exponents. For \(v\le4\) the real comparison
per face costs \(2^v\eta_v=O(j^{-p})\), also at integer size.
\end{proof}
